\subsection{Scope and Limitations of the Project}

\ChapterContent{
  The main focus of the project is the improvement of the inventory process of X University by implementing an inventory system that enhances the overall management, control, and efficiency of inventory-related operations within the institution. The system features include a centralized asset database, dynamic QR code generation, mobile-integrated scanning and auditing, and an automated depreciation and maintenance tracker. The project will utilize a mixed-methods research design, with the Office of Supply serving as the primary source of information and relevant documents, using interview questionnaires and observation as data-gathering instruments. The project framework combines the Extreme Programming methodology for system development and Business Process Reengineering for process improvement to ensure balanced focus on both system development and process enhancement, with implementation using the WAMP stack as the development environment. Lastly, the system quality will be evaluated by the Supply Office staff using the Functional Suitability and Interaction Capability characteristics of ISO 25010:2023, along with Time Efficiency, Accuracy, and Productivity criteria derived from Business Key Performance Indicators (KPI).\\

The project is limited to the inventory processes handled by the Office of Supply of X University and does not cover other administrative or academic units outside its operational scope. The system focuses primarily on asset monitoring, inventory tracking, and related reporting functions, excluding advanced financial accounting and procurement management modules. Implementation will be confined to the identified users and existing institutional infrastructure, which may limit integration with external systems and future organizational expansions. Data collection and evaluation will rely on the availability, participation, and feedback of Supply Office personnel during the development and testing phases. Additionally, the system evaluation will concentrate on selected ISO 25010:2023 characteristics and identified business KPIs, and therefore will not assess all possible software quality attributes or long-term operational impacts.

}